% Target language: Spanish (Spain), es-ES.
% Source: content/en-US.tex, checked against content/zh-CN.tex at upstream commit 0d064f6.
% Translation: AI-assisted; needs native-speaker review before publication.
\StoryParagraph{3}{El peso de un solo óvulo: 3,5 microgramos.}

\StoryParagraph{3}{El peso de cincuenta millones de dólares en efectivo: 2,5 toneladas.}

\StoryParagraph{3}{En una llamada desde Montage Laguna Beach, Jing Tian me pidió lo segundo y ofreció lo primero como garantía.}

\StoryParagraph{3}{Había ido allí para someterse a una extracción de óvulos. El bebé, que gestaría una mujer mediante gestación subrogada, estaba previsto para 2027. Si todo avanzaba lo bastante deprisa, aún podría nacer en el Año del Caballo. La idea había sido suya. Ya en la clínica, dijo: sin cincuenta millones, no hay extracción.}

\UnprintedPause



\StoryParagraph{0}{Montage Laguna Beach se alza sobre un acantilado frente al Pacífico. Reservamos la planta entera durante treinta días, por un millón de dólares, para proteger su intimidad. Se marchó al cabo de ocho.}

\StoryParagraph{0}{Durante los veintidós días restantes, el personal de limpieza siguió subiendo a diario, como siempre. Me lo contó mi asistente. Lo vio una vez, cuando fue a devolver la llave: una mujer mexicana empujaba un carro y entraba en una habitación tras otra, cambiaba las toallas, limpiaba las superficies, remetía las esquinas de las sábanas, salía y cerraba cada puerta a su espalda. Dos horas para hacer una planta entera.}

\StoryParagraph{0}{Mi asistente le preguntó por qué seguía haciéndolo si no se alojaba nadie.}

\StoryParagraph{0}{Ella respondió que, según el registro, sí había alguien.}

\StoryParagraph{0}{Un Versalles que había perdido a su reina.}

\StoryParagraph{0}{Entonces comprendí que a Jing Tian le gustaba aquello. A Jing Tian le gustaba que un lugar permaneciera vacío por ella. Que no ocurriera nada.}

\StoryParagraph{3}{Tan vacío como Montage Laguna Beach estaba un Gulfstream G700, detenido con aire incómodo en Van Nuys mientras su tripulación se presentaba cada día a trabajar con toda normalidad.}

\StoryParagraph{3}{Al octavo día hizo aquella llamada. Después, nadie volvió a verla.}

\NarrativeBreak



\StoryParagraph{0}{En 2007 ya había alguien quemando dinero por ella. Aquel año yo acababa de entrar en la Universidad de Pekín.}

\StoryParagraph{3}{Por supuesto, no era yo quien quemaba el dinero.}

\StoryParagraph{3}{QQ abrió una ventana emergente con ella como imagen de la marca. Fue en aquel viejo Lenovo IBM, con una resolución tan mala que todo se veía borroso. Guardé la imagen.}

\StoryParagraph{3}{Desde mi dormitorio de la Universidad de Pekín, cuyo alquiler anual era de mil yuanes, entré en internet, la encontré en Xiaonei y le envié una solicitud de amistad. Xiaonei obligaba a añadir un mensaje. Escribí y borré, borré y volví a escribir durante unos cuarenta minutos. Al final envié: «Hola».}

\StoryParagraph{3}{No aceptó la solicitud.}

\StoryParagraph{3}{Puede que ni siquiera hubiera usado Xiaonei.}

\StoryParagraph{3}{Durante los años siguientes, hiciera lo que hiciese, no dejé de pensar en una pregunta: ¿hasta dónde tendría que llegar para que me aceptara? Lo pensé al alcanzar mis primeros mil millones. Volví a pensarlo al llegar a los diez mil millones.}

\StoryParagraph{3}{Con el tiempo dejé de pensarlo. Renren había quebrado hacía años.}

\StoryParagraph{3}{Nunca envié una segunda solicitud.}

\StoryParagraph{3}{En 2011 hacía frío en Pekín. Compré un plumífero por ciento cincuenta yuanes. Recuerdo el precio porque dudé durante tres días.}

\StoryParagraph{3}{Ese mismo año protagonizó una película de ciento cincuenta millones de yuanes, \textit{The Warring States}. Sun Honglei, Francis Ng, Kim Hee-sun, Kiichi Nakai y Jiang Wu habían sido convocados como mero decorado, igual que Montage Laguna Beach y el Gulfstream G700 diecinueve años después. No sabían por qué estaban allí. Ya que habían venido, algo tendría que suceder. No sucedió nada.}

\StoryParagraph{3}{Después llegó \textit{The Great Wall}, con ciento cincuenta millones de dólares. Zhang Yimou. Andy Lau. Matt Damon. Después, Jackie Chan. La lista del universo de Jing Tian no paraba de crecer.}

\StoryParagraph{3}{Para entonces yo ya estaba en Estados Unidos: estudiaba, fundaba una empresa, cada vez tenía más trabajo. Lo único que nunca cambió fue que, cada vez que cambiaba de dispositivo, por muchas molestias que supusiera, transfería aquella imagen de QQ.}

\StoryParagraph{3}{Diecinueve años después, mi nombre se incorporó a la lista.}

\NarrativeBreak

\StoryParagraph{3}{Nos vimos por primera vez en el Rosewood Hong Kong. Siempre había pensado que estaría nervioso el día que la conociera.}

\StoryParagraph{3}{No lo estaba.}

\StoryParagraph{3}{Estaba más delgada que en las fotos, vestía de manera informal y descansaba recostada en un sofá. La miré y, de pronto, no pude recordar qué era exactamente lo que llevaba tanto tiempo esperando.}

\StoryParagraph{3}{La llevé a ver el plátano de Cattelan. Lo contempló durante un buen rato y después me preguntó: «¿Cuánto cuesta esto?»}

\StoryParagraph{3}{Seis millones doscientos mil dólares.}

\StoryParagraph{3}{Dijo: «¿Y qué pasa cuando se pudre?»}

\StoryParagraph{3}{No respondí. Ella dijo: «Déjalo. Vamos a ver una película».}

\StoryParagraph{3}{Dije: «Vale».}

\StoryParagraph{3}{Dijo que en la sala no podía haber nadie salvo nosotros dos.}

\StoryParagraph{3}{Dije: «Vale».}

\StoryParagraph{3}{Salí para hacer una llamada. Cuando regresé, seguía de pie delante del plátano.}

\StoryParagraph{3}{La película era \textit{Zootrópolis 2}.}

\StoryParagraph{3}{Una sala de proyección para dos. Ese tipo de escena suele estar reservado a los pases de censura del Departamento Central de Propaganda.}

\StoryParagraph{3}{Hice que mi asistente comprara todas las butacas del cine.}

\StoryParagraph{3}{Había leído artículos que denunciaban la falsificación de la taquilla: todas las entradas vendidas y solo dos personas en la sala. Ahora lo entendía. Tal vez no siempre fuera por la taquilla.}

\StoryParagraph{3}{Comprar todas las butacas no fue difícil. El problema eran las entradas ya vendidas. Mis asistentes y guardaespaldas llegaron temprano y esperaron en la puerta. Cada vez que aparecía alguien con una entrada, le entregaban dinero en efectivo.}

\StoryParagraph{3}{Llegó una pareja con su hijo. El niño ya estaba en la puerta, alzando un cubo de palomitas.}

\StoryParagraph{3}{Una asistente les tendió un fajo de billetes. El niño no entendía por qué no podían entrar. Los padres cogieron el dinero, dieron las gracias, no preguntaron nada y se marcharon.}

\StoryParagraph{3}{Aquel día aceptaron el dinero veintiséis personas.}

\StoryParagraph{3}{Después de la película, cambió su foto de perfil de WeChat por una de la agente Judy Hopps.}

\StoryParagraph{3}{Aquella noche le toqué las tetas con la mano. Ella me apartó.}

\StoryParagraph{3}{Pensé que no quería. Retiré la mano y estaba a punto de decir algo.}

\StoryParagraph{3}{No dijo nada. Se incorporó, sacó de su bolso un pequeño estuche y abrió la cremallera. Dentro había una hilera de objetos, como un diminuto juego de herramientas.}

\StoryParagraph{3}{Dijo: «Tus uñas me hacen daño».}

\StoryParagraph{3}{Miré. Siempre me las había cortado de cualquier manera con un cortaúñas. Los bordes quedaban afilados. Normalmente ni lo notaba.}

\StoryParagraph{3}{Dijo: «Ven».}

\StoryParagraph{3}{Me puso la mano sobre el regazo y me limó las uñas una a una.}

\StoryParagraph{3}{Primero usó la cara áspera. Cuando terminó con ella, cambió de cara y después volvió a cambiar. De vez en cuando levantaba un dedo hacia la luz, lo examinaba y seguía limando.}

\StoryParagraph{3}{Pregunté: «¿Cuánto vas a tardar?»}

\StoryParagraph{3}{Dijo: «Hasta que dejen de arañar a la gente».}

\StoryParagraph{3}{Diez dedos. Tardó casi una hora.}

\StoryParagraph{3}{Mientras trabajaba, me enseñó: esta cara sirve para dar forma, esta para pulir y, con la última, no hay que frotar de un lado a otro. Hay que pasarla en una sola dirección.}

\StoryParagraph{3}{Dije: «Vale».}

\StoryParagraph{3}{Cuando terminó, me dio la vuelta a la mano, la examinó y pasó los dedos por encima.}

\StoryParagraph{3}{Dijo: «Ya está».}

\StoryParagraph{3}{«Adelante. Toca donde quieras».}

\StoryParagraph{3}{En un momento dado dijo algo. Dijo: «Si algún día ya no estoy, nunca volverá a haber nadie que te pula las uñas».}

\StoryParagraph{3}{Sonreí. No dije nada.}

\StoryParagraph{0}{Tres meses después de conocerla, encargué que prepararan un informe. Sus proyectos a lo largo de los años, sus socios comerciales, sus participaciones empresariales, Jing Guoqing, Tian Xin'ai, todas las cuentas relacionadas, la vivienda alquilada en Pekín, la comprada en Shanghái y algunas cosas más. Tardaron una semana. Cuarenta y tantas páginas en total. También había asuntos de hombres y mujeres. Me parecieron lo menos importante.}

\StoryParagraph{0}{Más adelante me habló de algunas de las cosas que figuraban allí. Lo que me contó no coincidía del todo con lo que estaba escrito. No dije nada.}

\StoryParagraph{0}{Hubo cosas de las que nunca me habló. Yo nunca pregunté.}

\StoryParagraph{0}{Creí todo lo que me dijo.}

\StoryParagraph{0}{Guardé aquel informe. Aún lo conservo.}

\StoryParagraph{3}{Pero no creí ni una sola palabra.}

\NarrativeBreak

\StoryParagraph{3}{Dijo: «Hong Kong está abarrotado. Me reconocerán. Vayámonos a algún sitio».}

\StoryParagraph{3}{Dije: «Vale».}

\StoryParagraph{3}{Señaló Tenerife en el mapa.}

\StoryParagraph{3}{Dije: «Vale».}

\StoryParagraph{3}{Salí para hacer una llamada. Solo al regresar supe que aquel lugar estaba en el Atlántico, frente a la costa noroccidental de África. Al mediodía siguiente partimos los treinta. Un A330 nos esperaba en el aeropuerto.}

\StoryParagraph{3}{En el dormitorio del avión me preguntó: «¿Te gusta cómo suena mi nombre?»}

\StoryParagraph{3}{Dije: «Sí».}

\StoryParagraph{3}{Dijo: «Es el apellido de mi padre unido al de mi madre. ¿No es una coincidencia curiosa?»}

\StoryParagraph{3}{Dije: «Lo recordé desde el primer día».}

\StoryParagraph{3}{Me rodeó con los brazos y dijo: «Ya no me llames Jing Tian. Llámame mamá».}

\StoryParagraph{3}{«Vale, mamá».}

\StoryParagraph{3}{Después de aterrizar, el capitán se acercó y me dijo: «Jefe, al principio la isla no nos dejaba aterrizar. Nunca habían visto un avión privado tan grande».}

\StoryParagraph{3}{En la isla regían las mismas normas: no podía haber nadie cerca de nosotros. Mis asistentes se ocuparon. Nadie preguntó por qué. Esta vez yo no pregunté cuántos eran.}

\StoryParagraph{3}{El Ritz-Carlton daba al Atlántico y sus ventanales eran demasiado grandes. Ella no podía dormir si quedaba el menor rastro de luz. Los asistentes cubrieron toda la habitación con cinta adhesiva negra. Fuera resplandecía el día; dentro reinaba una oscuridad de celda de aislamiento. La cinta se despegaba con el sol y había que volver a colocarla a diario.}

\StoryParagraph{3}{Había once asistentes en total. Llegaban cada mañana a las seis y media y no se marchaban hasta después de que nos durmiéramos. Dos chicas se ocupaban de la cinta: una cubría los marcos de las ventanas; la otra, los pilotos del aire acondicionado y los enchufes. Al terminar, lo comprobaban todo tres veces. Compraron treinta rollos de cinta y luego otros veinte. Descubrieron que la que colocaban después de las tres de la tarde se despegaba con más facilidad porque a esa hora el cristal estaba más caliente, así que pasaron a hacerlo por la mañana. Nadie se lo enseñó. Lo averiguaron ellas solas. Lo comentaban entre sí como si fuese un oficio.}

\StoryParagraph{3}{Ella sabía cómo se llamaban aquellas dos chicas. Yo no.}

\StoryParagraph{3}{Nos detuvimos junto a la carretera del Parque Nacional del Teide para contemplar la puesta de sol. Los asistentes bajaron rosas del coche y abrieron champán.}

\StoryParagraph{3}{Mamá contempló el mar de nubes y dijo: «Qué bonito. Pero no tanto como la puesta de sol de Delingha, en Qinghai».}

\StoryParagraph{3}{Yo soy de Qinghai.}

\StoryParagraph{3}{No lo dije.}

\StoryParagraph{3}{Por la noche contemplamos las estrellas desde la cima. Llevábamos abrigos acolchados y aun así teníamos frío.}

\StoryParagraph{3}{Mamá dijo: «Pide un deseo».}

\StoryParagraph{3}{Pregunté: «¿Qué has pedido?»}

\StoryParagraph{3}{Mamá dijo: «No te lo voy a contar».}

\StoryParagraph{3}{Mamá dijo: «Es precioso. Pero hace demasiado frío. ¿Nos vamos a las Maldivas?»}

\StoryParagraph{3}{Dije: «De acuerdo».}

\StoryParagraph{3}{En Tenerife había quince grados bajo cero y en las Maldivas, veintisiete sobre cero. Nueve mil quinientos kilómetros en línea recta. Once horas de vuelo. Pero el aeropuerto de las Maldivas era demasiado pequeño, así que a mitad de camino tuvimos que cambiar nuestro avión de fuselaje ancho por otro más pequeño.}

\StoryParagraph{3}{Mamá dijo: «Los aviones pequeños hacen demasiado ruido. Los grandes son mejores».}

\StoryParagraph{3}{En el dormitorio del avión rumbo a las Maldivas, mamá me contó que su deseo había sido tener un hijo conmigo. Tenía que ser mediante gestación subrogada.}

\StoryParagraph{3}{Nunca le pregunté por qué. Pensé que tenía miedo al dolor.}

\StoryParagraph{3}{En pleno vuelo, mamá se incorporó de pronto y buscó el teléfono. Se conectó a Starlink y dijo: «Qué desastre».}

\StoryParagraph{3}{Pregunté: «¿Qué ocurre?»}

\StoryParagraph{3}{Dijo: «Hay alguien que me está amenazando. Lleva tres días insistiendo y no le he contestado».}

\StoryParagraph{3}{Dije: «¿Cuánto quiere?»}

\StoryParagraph{3}{Dijo: «No hace falta que te metas».}

\StoryParagraph{3}{Estuvo escribiendo durante mucho tiempo: borraba y cambiaba, cambiaba y borraba. Al terminar, dejó el teléfono boca abajo sobre el regazo. Al cabo de un rato volvió a darle la vuelta para mirar.}

\StoryParagraph{3}{Le pregunté quién era.}

\StoryParagraph{0}{Dijo: «Alguien que trabajaba conmigo. En finanzas. Entró en prisión. Cargó con la culpa».}

\StoryParagraph{3}{«Ahora ha salido, no tiene adónde ir e insiste en recurrir a mí».}

\StoryParagraph{3}{Dije: «Entonces llama a la policía y haz que lo detengan».}

\StoryParagraph{3}{Me miró. «No lo entiendes».}

\StoryParagraph{3}{Dije: «¿Le has pagado?»}

\StoryParagraph{3}{Dijo: «Sí».}

\StoryParagraph{3}{«¿Cuánto?»}

\StoryParagraph{3}{«Ochenta mil».}

\StoryParagraph{3}{Dije: «¿Has tardado tres días en decidir eso?»}

\StoryParagraph{3}{No dijo nada.}

\StoryParagraph{3}{Apenas habló durante el resto del día. No se durmió hasta que aterrizamos.}

\StoryParagraph{0}{La maleta de mamá fue lo primero que llegó a Soneva Jani. El equipaje de treinta personas no cabía en un solo avión pequeño, así que lo repartieron entre tres. Mamá solo llevaba una maleta. Aquella maleta tuvo un avión para ella sola. Dijo que no la pusieran junto a la de nadie más.}

\StoryParagraph{0}{Soneva Jani era famoso por su lema «no news, no shoes». Mamá se quedó tranquila.}

\StoryParagraph{3}{Entre mayordomos del hotel, guardaespaldas y asistentes sumaban veintiséis personas; además, dos operadores de dron nos seguían a distancia. Los asistentes empezaron de nuevo a taparlo todo con cinta: el piloto del aire acondicionado, las luces de los enchufes, una y otra vez.}

\StoryParagraph{0}{Por la noche intercambiamos historias sobre nuestras exparejas como quien intercambia contraseñas, como una dinastía que coteja los nombres de sus eras. Mamá no toleraba la menor discrepancia. Pero tenía una palabra prohibida: no se podía mencionar al deportista.}

\StoryParagraph{0}{No lo mencioné, aunque sabía que el tatuaje de la espalda de mamá se lo había hecho por él.}

\StoryParagraph{3}{Quizá era demasiado grande para borrarlo.}

\StoryParagraph{0}{Dijo: «Ahora somos una familia. A mis padres no les resultó fácil criarme. Tienes que darles el pago nupcial».}

\StoryParagraph{0}{Dije: «Vale».}

\StoryParagraph{0}{Me dio dos nombres, Jing Guoqing y Tian Xin'ai, y dos números de cuenta.}

\StoryParagraph{0}{Treinta millones. En renminbis.}

\StoryParagraph{0}{Hice la transferencia. Había visto por primera vez aquellos dos nombres en las cuarenta y tantas páginas.}

\StoryParagraph{3}{Al cabo de un rato me tendió el teléfono. Tian Xin'ai había respondido con una sola palabra: «Recibido».}

\StoryParagraph{0}{Una vez mamá dijo: «¿Sabes qué pienso cuando te pido estas cosas?»}

\StoryParagraph{0}{Dije: «¿Qué?»}

\StoryParagraph{0}{Dijo: «Me pregunto si algún día dirás que no».}

\StoryParagraph{0}{Dije: «No lo haré».}

\StoryParagraph{0}{Dijo: «Entonces no tiene gracia».}

\StoryParagraph{3}{Pensé: ¿qué podría haber que yo no fuese capaz de hacer?}

\StoryParagraph{0}{Cuando me aburrí demasiado, me puse a trastear con la imagen de QQ de 2007 para pasarla a un dispositivo nuevo. Mamá se inclinó hacia mí. «¿Por qué guardas a escondidas la foto de otra mujer? ¿Quién es?»}

\StoryParagraph{0}{Guardé silencio. No dije nada.}

\StoryParagraph{0}{La imagen estaba demasiado borrosa, así que se marchó a hacer otra cosa.}

\StoryParagraph{0}{Aquella noche borré la imagen.}

\StoryParagraph{3}{Al día siguiente la recuperé de un teléfono antiguo.}

\StoryParagraph{0}{Mamá tampoco dejaba de deslizar el dedo por la pantalla del teléfono. Le pregunté qué estaba mirando.}

\StoryParagraph{0}{Dijo: «Nada. Toda esa mierda de pelearse por encabezar los créditos».}

\StoryParagraph{0}{Pelearse por los créditos: en una película solo puede haber una estrella cuyo nombre figure primero. Incluso entre el protagonista y la protagonista, solo uno puede aparecer en primer lugar.}

\StoryParagraph{0}{Dije: «Qué aburrimiento».}

\StoryParagraph{0}{Dijo: «Sí. Y lo peor es que a veces se forman parejas de rodaje. Pasas cuarenta o cincuenta días filmando, casi siempre frente al mismo hombre. Por la noche te acuestas con él y durante el día os peleáis por quién encabeza los créditos».}

\StoryParagraph{0}{Dije: «Eso es inhumano».}

\StoryParagraph{0}{Dijo: «Sí».}

\StoryParagraph{0}{«Pero cuando termina el rodaje, da igual cuántas veces os hayáis acostado: fingís que nunca ocurrió nada».}

\StoryParagraph{3}{Después siguió leyendo durante otra media hora.}

\StoryParagraph{0}{Año nuevo. El 2 de enero de 2026, en Soneva Jani, le pedí matrimonio a mamá.}

\StoryParagraph{0}{Lo había decidido la mañana del 1 de enero. El anillo estaba en Hong Kong. En la isla no había ningún sitio donde comprar uno, así que una asistente tomó un hidroavión hasta Malé. Regresó aquella tarde con el anillo dentro de un sobre del hotel.}

\StoryParagraph{0}{Mamá dijo: «El diamante es demasiado pequeño».}

\StoryParagraph{0}{Mamá dijo: «No pasa nada. Ya me lo compensarás más adelante con uno más grande».}

\StoryParagraph{0}{«Y otra cosa: ¿puedes venir a Pekín?»}

\StoryParagraph{0}{Guardé silencio.}

\StoryParagraph{0}{Pensó que no quería ir.}

\StoryParagraph{0}{Mamá dijo: «¿Por qué eres tan bueno conmigo?»}

\StoryParagraph{0}{Dije: «Porque me gustas».}

\StoryParagraph{0}{Dijo: «¿También eras así con otras personas?»}

\StoryParagraph{0}{Dije: «No».}

\StoryParagraph{3}{Sonrió y no dijo nada.}

\StoryParagraph{0}{Una mañana me desperté temprano. Ella no estaba en la cama.}

\StoryParagraph{0}{Salí y la encontré sentada en la terraza, frente al mar. El mar entero era negro. No había ninguna luz encendida. Aún no había amanecido.}

\StoryParagraph{0}{Le pregunté por qué se había levantado tan pronto.}

\StoryParagraph{0}{Dijo: «Estoy contando los días».}

\StoryParagraph{0}{Dije: «¿Contando qué?»}

\StoryParagraph{0}{Dijo: «Vuelve a dormir».}

\StoryParagraph{0}{Me senté a su lado durante un rato, sin saber qué decir. Después volví a la cama.}

\StoryParagraph{3}{Cuando desperté de nuevo, ya se había maquillado.}

\StoryParagraph{0}{Durante los catorce días que pasamos en las Maldivas, ni una sola persona la reconoció.}

\StoryParagraph{0}{La decimotercera noche, en el restaurante, se quitó las gafas de sol. Al cabo de un rato se quitó también el sombrero y lo dejó sobre la mesa. Permaneció así durante toda la cena. En la mesa de al lado había una pareja de ancianos alemanes. El camarero era de Sri Lanka. Nadie volvió la cabeza para mirarla.}

\StoryParagraph{0}{Después de cenar volvió a ponerse el sombrero.}

\StoryParagraph{0}{Mamá dijo: «Volvamos».}

\StoryParagraph{0}{Dije: «¿No se suponía que nos quedaríamos más tiempo?»}

\StoryParagraph{3}{Dijo: «Hay demasiado silencio».}

\UnprintedPause

\StoryParagraph{0}{De vuelta en Hong Kong, mamá conoció a mis padres.}

\StoryParagraph{0}{Aquella noche dijo que su familia tenía que preparar el encuentro entre las dos familias y la boda. Yo debía enviar todo lo posible. Los destinatarios seguían siendo Jing Guoqing y Tian Xin'ai.}

\StoryParagraph{0}{No pregunté por qué.}

\StoryParagraph{0}{Dije: «El límite diario de cada tarjeta es de un millón».}

\StoryParagraph{0}{Mamá dijo: «Demasiado lento».}

\StoryParagraph{0}{«Usa cinco tarjetas».}

\StoryParagraph{0}{Cinco tarjetas, dos cuentas, alternando entre ellas. Al llegar a la tercera tarjeta, dudé. Aun así, completé todas las transferencias.}

\StoryParagraph{3}{Mamá dijo: «Demasiado ineficiente».}

\StoryParagraph{0}{Dije: «Traigamos a nuestros padres a Hong Kong para pasar juntos el Año Nuevo».}

\StoryParagraph{0}{Mamá dijo: «Este año ya está decidido que lo pasaremos en Hainan. El próximo».}

\StoryParagraph{0}{Durante el Año Nuevo chino, los fuegos artificiales sobre el puerto Victoria congregaron a un millón de personas. Ella no estaba entre ellas.}

\StoryParagraph{0}{A las tres de la madrugada llevé a mis padres a casa.}

\StoryParagraph{0}{Las calles estaban vacías. No había ningún otro coche esperando ante el semáforo en rojo. Aun así, me detuve.}

\StoryParagraph{0}{Mamá dijo que, por el momento, sería mejor que no nos viéramos. Estaba a punto de ir a Los Ángeles para la extracción de óvulos. Antes tenía asuntos que resolver.}

\StoryParagraph{0}{Mamá volvió a preguntar: «¿Puedes venir a Pekín?»}

\StoryParagraph{0}{Guardé silencio.}

\StoryParagraph{3}{Nunca dijo qué había en Pekín. Yo nunca pregunté.}

\StoryParagraph{0}{Mamá dijo que teníamos que comprar una vivienda en Hong Kong y ponerla a su nombre.}

\StoryParagraph{0}{Dije: «De acuerdo».}

\StoryParagraph{3}{Tampoco nos vimos el día de San Valentín.}

\StoryParagraph{0}{Una vez, por teléfono, dijo: «Esta mañana me han sacado ocho tubos de sangre».}

\StoryParagraph{0}{Dije: «¿Tantos?»}

\StoryParagraph{0}{Dijo: «No ha sido nada».}

\StoryParagraph{0}{Dije: «Ha debido de ser duro».}

\StoryParagraph{0}{Después hablamos de otra cosa.}

\StoryParagraph{0}{Solo más tarde supe que, para aquel tipo de extracción de sangre, había que estar en ayunas. Tenía que hacerse el segundo o el tercer día del ciclo menstrual; no podía equivocarse de fecha. Probablemente salió de casa hacia las seis de la mañana. Invierno en Pekín. Aún de noche. Después llegaron cada vez más pruebas. Lo supe porque apareció una partida nueva en el resumen mensual de pagos de mi asistente. La primera vez fue el año anterior. No era mucho, unos pocos miles de yuanes. A partir de entonces aparecía todos los meses, a veces dos o tres cargos. Los aprobaba deprisa. Cantidades así no exigían que revisara los detalles. Solo más adelante supe que se trataba de análisis de sangre, un perfil de seis hormonas, ecografías y AMH. Había que averiguar cuánta reserva ovárica le quedaba, cuántos óvulos podrían extraerle y si merecía la pena seguir adelante.}

\StoryParagraph{0}{Todo se hizo en Pekín.}

\StoryParagraph{3}{No sé en qué hospital. No sé si fue sola o si alguien la acompañó.}

\StoryParagraph{0}{El 25 de febrero aterrizó en Hong Kong.}

\StoryParagraph{0}{Dijo que no iría a Los Ángeles hasta ver antes la vivienda de Hong Kong.}

\StoryParagraph{0}{Le envié la docena larga de viviendas que habíamos examinado durante los seis meses anteriores. Todas le parecieron feas. Yo no había visitado ni una. Mis asistentes habían pasado seis meses recorriéndolas una por una. Ellos me las enviaban a mí; yo se las enviaba a ella. Una estaba en Mid-Levels. En el vídeo se veía el puerto Victoria. Lo habían grabado al anochecer. Por un instante la cámara se movió y captó en la pared la sombra de la persona que filmaba. Contemplé aquella sombra y pensé: ¿quién es esta persona? ¿Cuánto ha caminado hoy? ¿Cuántas viviendas ha visto? ¿Sabe para quién están enseñando esta?}

\StoryParagraph{3}{Quedaba demasiado poco tiempo. Aun así, subió al G700 rumbo a Los Ángeles y voló a Montage Laguna Beach el 28 de febrero.}

\NarrativeBreak

\StoryParagraph{0}{Al octavo día llamó mamá. Para ella era por la tarde; para mí, plena madrugada.}

\StoryParagraph{0}{Dijo: «Cincuenta millones. De dólares».}

\StoryParagraph{0}{Dije: «Déjame pensarlo».}

\StoryParagraph{0}{No dijo nada. Oía el océano al otro lado.}

\StoryParagraph{0}{Al cabo de un rato dijo: «Vale».}

\StoryParagraph{0}{Colgó.}

\StoryParagraph{3}{Desde que la conocí, jamás le había dicho «déjame pensarlo».}

\StoryParagraph{0}{Usé Claude Code para consultar las API y calcular todos mis activos en efectivo. La conclusión fue que aquello no tendría ningún efecto sobre mí. Volví a comprobarlo. Escribí la pregunta en Claude Code. Claude dijo: «No le des los cincuenta millones de dólares».}

\StoryParagraph{3}{Pregunté: «¿Ya no me quiere?»}

\StoryParagraph{3}{Claude dijo: «En lo que respecta al amor, ni me importa ni lo entiendo. Pero no puedes darle esos cincuenta millones de dólares».}

\StoryParagraph{3}{Nunca me había negado a darle nada. No por generosidad. Por miedo. Miedo a la respuesta que llegaría un segundo después.}

\StoryParagraph{0}{Y ahora alguien lo había dicho por mí.}

\StoryParagraph{0}{Contemplé aquella línea durante mucho tiempo. Lo que me aterrorizaba no era que Claude fuese inteligente. Era que Claude no había dudado en absoluto.}

\StoryParagraph{0}{No tenía dudas.}

\StoryParagraph{0}{Pensé que aquello debía de ser una forma superior de existencia. Los seres humanos no podían hacerlo.}

\StoryParagraph{3}{¿Ya no me quería?}

\StoryParagraph{3}{Al día siguiente volvió a llamar.}

\StoryParagraph{3}{Mamá dijo: «El océano es inmenso».}

\StoryParagraph{3}{Dije: «¿Ah, sí?»}

\StoryParagraph{3}{Dijo: «Esta planta es demasiado grande. Hay eco. Esta mañana ha venido alguien a limpiar. Me he escondido en el cuarto de baño y he esperado a que se marchara».}

\StoryParagraph{3}{Dije: «Puedo decirles que no vengan».}

\StoryParagraph{3}{Dijo: «No hace falta».}

\StoryParagraph{3}{Dijo: «La clínica ha preguntado cuándo llegará el padre del bebé».}

\StoryParagraph{3}{Dije: «Puedo enviar a alguien».}

\StoryParagraph{3}{Dijo: «No hace falta».}

\StoryParagraph{3}{Entonces dijo: «¿Ya lo has decidido?»}

\StoryParagraph{3}{No hablé. No porque estuviera pensando. No estaba pensando en nada. Por teléfono, en su extremo estaba el océano y en el mío, el puerto Victoria. Ninguno de los dos hablaba. Conté. Once segundos.}

\StoryParagraph{3}{Jing Tian dijo: «Está bien».}

\StoryParagraph{3}{Jing Tian dijo: «Lo entiendo».}

\StoryParagraph{3}{Colgó.}

\StoryParagraph{3}{He vuelto innumerables veces a aquellos dos días, intentando recordar qué dijo aparte de lo de los cincuenta millones. El océano es inmenso. Esta planta es demasiado grande. Se escondió en el cuarto de baño. Eso fue todo. Recuerdo cada palabra. Lo que no puedo recordar es su voz, el tono con el que dijo: «Lo entiendo».}

\StoryParagraph{0}{No puedo recordarlo. Al final dejé de intentarlo.}

\StoryParagraph{0}{Dejé el teléfono boca abajo sobre la mesilla. Después hice algo que todavía no comprendo: me puse una mano sobre el pecho y conté.}

\StoryParagraph{0}{El corazón latía con regularidad.}

\StoryParagraph{0}{Volví a contar. Seguía latiendo con regularidad. Algo más de sesenta pulsaciones por minuto, como siempre. Permanecí allí sentado, con la mano apretada contra el pecho, como un médico que examinara a otra persona.}

\StoryParagraph{3}{Pensé: así que era esto.}

\StoryParagraph{0}{No hice nada durante el resto del día. Me quedé en la habitación hasta que oscureció y volvió a amanecer, mientras contemplaba cómo el puerto Victoria pasaba del azul al gris y después al negro. A las seis de la mañana, las torres financieras de la otra orilla empezaron a iluminarse una tras otra. Una asistente llamó una vez a la puerta para preguntar qué debía preparar para cenar. Dije que no hacía falta.}

\StoryParagraph{3}{No lloré. Seguí esperando a que me diera por llorar.}

\NarrativeBreak

\StoryParagraph{0}{Más tarde, cuando mi asistente ajustó las cuentas conmigo, una de las partidas era una indemnización al hotel de Tenerife. Después de marcharnos, siguieron colocando la cinta durante varios días. Nadie les dijo a las chicas que podían parar. El hotel afirmó que la habían pegado durante demasiado tiempo, demasiadas veces. El residuo del marco de la ventana no salía. Había que sustituir el cristal entero.}

\StoryParagraph{0}{Dije: «Aprobado».}

\StoryParagraph{3}{Aquel cristal daba al Atlántico.}

\StoryParagraph{0}{El combustible del avión que figuraba en la factura costaba más que toda la planta del Montage. Pregunté por ello. Mi asistente dijo que el A330 estaba diseñado para doscientas personas. Con solo treinta a bordo, pesaba demasiado poco y el equilibrio para el despegue no era correcto. Por eso, cada vez tenían que añadir combustible hasta que pesara lo suficiente. Al llegar al destino, el combustible sobrante no podía quedarse a bordo. Había que desecharlo.}

\StoryParagraph{0}{Dije: «¿Todas las veces?»}

\StoryParagraph{3}{Él dijo: «Todas las veces».}

\StoryParagraph{0}{También había una caja enviada desde la clínica de Los Ángeles. Ocho días de inyecciones de estimulación ovárica. Tenían que conservarse refrigeradas. No habían abierto ni una.}

\StoryParagraph{3}{También estaba el depósito de la agencia de gestación subrogada: doscientos mil dólares, no reembolsables. Habían elegido a la gestante en enero. Era de California, tenía treinta y cuatro años y había dado a luz a dos hijos. Su fotografía estaba en el expediente. Nunca la miré. Llevaba preparada desde enero.}

\StoryParagraph{3}{Después de que Jing Tian se marchara, saqué la fotografía de 2007 y la miré una vez. Habían pasado diecinueve años. Los píxeles de todos los dispositivos no dejaban de mejorar, mientras que los de aquella imagen parecían empeorar. Su rostro estaba borroso. Amplié la foto y la observé durante mucho tiempo, buscando alguna diferencia entre la mujer de entonces y la de después. No encontré ni una. La fotografía nunca había sido lo bastante nítida para verla.}

\StoryParagraph{0}{Un día me corté las uñas. Al terminar pasé un dedo por encima. Estaban afiladas.}

\StoryParagraph{0}{Recordé que aquel objeto seguía en el cajón, el que ella había dejado.}

\StoryParagraph{0}{Lo saqué y lo pasé dos veces por las uñas. No parecía correcto. Ella había dicho que había que moverlo en una sola dirección. Había olvidado cuál.}

\StoryParagraph{0}{Volví a guardarlo.}

\StoryParagraph{3}{Mis uñas todavía arañan a la gente.}

\StoryParagraph{0}{Una mañana apareció un recordatorio en mi teléfono. Mi asistente lo había programado el año anterior al organizar el calendario: tres meses antes de la fecha prevista del parto, preparar la habitación.}

\StoryParagraph{3}{Deslicé el dedo para quitarlo.}

\StoryParagraph{0}{Desde que se marchó, no he dejado de pensar en una pregunta.}

\StoryParagraph{0}{Si aquel día le hubiera dado el dinero, ¿se habría quedado?}

\StoryParagraph{0}{Lo pensé durante mucho tiempo.}

\StoryParagraph{3}{Más tarde saqué el informe y volví a leerlo. Cuarenta y tantas páginas.}

\StoryParagraph{0}{Entonces recordé los fuegos artificiales sobre el puerto Victoria la noche del Año Nuevo chino, vistos desde la última planta.}

\StoryParagraph{0}{Un millón de personas llegaron como una marea. Terminaron los fuegos artificiales. Un millón de personas se retiraron.}

\StoryParagraph{0}{Un millón de personas.}

\StoryParagraph{0}{Todos los años es igual.}

\StoryParagraph{0}{Hong Kong a las tres de la madrugada:}

\StoryParagraph{3}{No ocurrió nada.}
